\documentclass[12pt,a4paper]{article}

\usepackage{fontspec}
\setmainfont{Libertinus Serif}
\usepackage[greek.ancient,spanish]{babel}
\babelfont[greek]{rm}{Libertinus Serif}
\usepackage[a4paper,margin=3.2cm]{geometry}
\usepackage{microtype}
\usepackage[hidelinks]{hyperref}

% Griego politónico escrito directamente en UTF-8:
% \gr{...} para palabras sueltas; bloques con \foreignlanguage{greek}{...}
\newcommand{\gr}[1]{\foreignlanguage{greek}{#1}}
\newcommand{\stanzabreak}{\par\vspace{0.6\baselineskip}}

\setlength{\parindent}{1.25em}
\setlength{\parskip}{0.35em}
\raggedbottom


\hypersetup{pdftitle={Contraportada (borrador): Leer la Odisea en tiempos iletrados}}

\pagestyle{empty}

\begin{document}

\section*{Contraportada (borrador)}

Cuando Christopher Nolan estrenó su \emph{Odisea}, las librerías se llenaron de Homeros. Comprarlos es fácil. Leerlos, en estos tiempos iletrados, es otra historia.

El autor de este libro leyó la \emph{Odisea} a los quince años, declamando la vieja prosa de Segalá a grito limpio mientras perseguía a su primo por las playas del Mar Menor. Medio siglo después ha vuelto al poema con la traducción en verso de Juan Manuel Rodríguez Tobal y se ha encontrado un texto más brutal, más moderno y más vivo de lo que recordaba.

Este libro es la síntesis de esas dos lecturas y de todo lo que pasó entre una y otra. Es un comentario de la \emph{Odisea}, episodio a episodio, que no pide al lector más que curiosidad. Son también unas memorias: las de un lector al que el poema le ha ido marcando el rumbo. Y es un ajuste de cuentas en verso con Ulises: si le quitamos los dioses, el héroe se queda sin coartada, y lo que queda a la vista es un hombre.

No hace falta haber leído la \emph{Odisea} para leer este libro. Lo difícil, después, será no leerla.

\end{document}
